\section{System Design: Proof-System Architecture}
\label{sec:design}

This section presents the full proof-system design of the Ziren WHIR
pipeline.  Where prior sections (\S\S\ref{sec:zerocheck-phase2a}--\ref{sec:whir-late-binding})
present each construction in isolation, here we describe how they
compose into a single end-to-end prover/verifier flow, with explicit
references to the implementation files.

\subsection{Proof-System Overview}

A Ziren shard proof attests that one shard of a MIPS program
executed correctly.  The proof contains:
\begin{itemize}
  \item A polynomial commitment to the per-chip main traces.
  \item Sumcheck-based proofs of transition-constraint and lookup
        soundness.
  \item Per-row multilinear-extension (MLE) openings of the main
        trace at GKR-derived points, bound back to the main commit
        via a late-binding wrapper.
  \item Cross-binding spot-check evidence linking the FRI commit
        and the WHIR late-binding commit to the same trace
        (production target).
\end{itemize}

The pipeline replaces three FRI-pipeline components with multilinear
equivalents:
\begin{enumerate}
  \item \textbf{Quotient identity} ($Q(\zeta) \cdot Z_H(\zeta) =
        C(\zeta)$) is replaced by zerocheck:
        $\sum_{\mathbf{b} \in \{0,1\}^m} \mathsf{eq}(\mathbf{r},
        \mathbf{b}) \cdot C(\mathbf{b}) = 0$, verified via sumcheck.
  \item \textbf{Permutation-trace lookup arguments} are replaced by
        LogUp-GKR with a per-layer sumcheck reduction.
  \item \textbf{Univariate FRI opening at $\zeta$} is augmented (not
        replaced --- the dual-commit approach keeps both for now)
        by WHIR multilinear opening at the GKR-derived row coordinates
        $\mathbf{r}_{\mathrm{row}}$.
\end{enumerate}

\subsection{Data Layout: Jagged Polynomial Commitment}
\label{sec:design-jagged}

Each shard contains $N$ chip traces of differing heights $h_c$ and
widths $w_c$, $c = 0, \ldots, N-1$.  Naively committing each chip
separately requires $N$ merkle trees and $N$ PCS instances.  We
adopt the Jagged construction~\cite{jagged} to amortise this:

\begin{construction}[Jagged trace packing]
\label{cons:jagged-pack}
Let $\mathbf{T}_0, \ldots, \mathbf{T}_{N-1}$ be per-chip trace
matrices.  The packed dense vector $\mathbf{q}$ is the column-major
concatenation of all per-chip columns:
\[
  \mathbf{q} \;=\;
  \mathbf{T}_0[:,0] \,\Vert\, \mathbf{T}_0[:,1] \,\Vert\, \ldots
   \,\Vert\, \mathbf{T}_{N-1}[:,w_{N-1}-1].
\]
Cumulative offsets $t_k$ satisfy $t_0 = 0$, $t_k = t_{k-1} +
h_{c(k)}$ where $k$ ranges over global column indices.  Padding
brings $|\mathbf{q}|$ to the next power of two $2^n$.
\end{construction}

The sparse $(\mathrm{row}, \mathrm{col})$-indexed polynomial $p(\mathbf{z}_r,
\mathbf{z}_c)$ relates to $\mathbf{q}$'s MLE $\widetilde{\mathbf{q}}$ via
\[
  p(\mathbf{z}_r, \mathbf{z}_c) \;=\;
  \sum_{j=0}^{|\mathbf{q}|-1}
    \widetilde{\mathbf{q}}(j) \cdot
    \mathsf{eq}(\mathrm{row}(j), \mathbf{z}_r) \cdot
    \mathsf{eq}(\mathrm{col}(j), \mathbf{z}_c)
\]
where $\mathrm{row}(j) = j - t_{\mathrm{col}(j)}$ and
$\mathrm{col}(j) = \min_k\{t_k > j\}$.  Recovering per-chip
evaluations from $\widetilde{\mathbf{q}}$ requires a sumcheck
argument that we describe in \S\ref{sec:design-jagged-sumcheck}.

\paragraph{Implementation.}
Layout is in \texttt{crates/stark/src/jagged.rs::pack\_traces\_jagged};
the resulting \texttt{JaggedPacking\{dense\_values, chip\_infos,
offsets, total\_values, log\_dense\_size\}} carries everything the
verifier needs to reconstruct $(\mathrm{row}(j), \mathrm{col}(j))$.
The chip-info hash (\texttt{hash\_chip\_infos}) is observed into the
challenger before commit to bind chip dimensions to the transcript.

\subsection{Polynomial Commitment: WHIR}
\label{sec:design-whir}

WHIR~\cite{whir2024} is a multilinear PCS based on Reed--Solomon
proximity testing.  A WHIR commitment to $\mathbf{q} \in
\mathbb{F}^{2^n}$ encodes $\mathbf{q}$'s coefficients into an RS
codeword on a coset of size $2^{n+r}$ ($r$ = log-inv-rate, set to $4$
in our deployment for $80$-bit security) and merkle-commits the row
hashes of the codeword matrix.

The WHIR open phase combines:
\begin{enumerate}
  \item Sumcheck folding rounds (one per WHIR round, each round
        folds $f$ variables for folding-factor $f$).
  \item STIR queries at $K$ random codeword positions, with
        per-round merkle paths.
  \item Out-of-domain (OOD) sample evaluations to bind the
        polynomial against the codeword.
  \item Per-round proof-of-work (PoW) grinding.
  \item A final-round in-the-clear polynomial of degree
        $2^{\mathrm{final\_sumcheck\_rounds}}$.
\end{enumerate}

\paragraph{Folding-factor scaling.}  Our jagged shards have
$2^n \in \{2^{22}, \ldots, 2^{27}\}$, but KoalaBear's two-adicity
caps at $24$.  We scale the folding factor as
$f = \max(4, n + r - 24)$ to keep the folded codeword in the
two-adic subgroup.  For $n = 27$: $f = 7$, giving
$\lceil n / f \rceil = 4$ rounds.

\paragraph{Implementation.}
WHIR PCS configuration is in
\texttt{crates/stark/src/whir\_config.rs}; the late-binding
adapter is in \texttt{crates/stark/src/whir\_late\_binding.rs}
(see \S\ref{sec:design-late-binding}).

\subsection{Late-Binding WHIR}
\label{sec:design-late-binding}

\paragraph{Motivation.}
The upstream Plonky3 \texttt{MultilinearPcs::commit} adapter takes
opening points at commit time, absorbing them into Fiat-Shamir
\emph{before} the column-batching challenge $\alpha$.  Our protocol
derives the row coordinates $\mathbf{r}_{\mathrm{row}}$ via GKR sumcheck
\emph{after} the commit, which the upstream adapter cannot accommodate.

\begin{construction}[Late-binding WHIR]
\label{cons:late-binding}
Phase order:
\begin{enumerate}
  \item \emph{Domain separator}: observe protocol parameters into the
        challenger.
  \item \emph{Commit phase}: build the RS codeword merkle and
        observe its root.  Sample OOD points; observe their
        evaluations into the challenger.  At this point the
        EqStatement contains only OOD constraints.
  \item \emph{(Caller-driven) post-commit point derivation}: GKR/
        zerocheck runs against the now-advanced challenger,
        producing $\mathbf{r}_{\mathrm{row}}$.
  \item \emph{Late-binding registration}: for each derived point
        $\mathbf{p}$ and value $v$, the prover invokes
        $\mathtt{add\_late\_claim}(\mathbf{p}, v)$, which evaluates
        $\widetilde{\mathbf{q}}(\mathbf{p}) = v$ (recorded into the
        EqStatement) and observes $v$ into the challenger.  The
        point is \emph{not} re-observed because it was derived from
        the post-commit challenger.
  \item \emph{Open phase}: run the standard WHIR open over the
        accumulated EqStatement (OOD claims + late claims).
\end{enumerate}
\end{construction}

\paragraph{Soundness.}  Identical to upstream WHIR.  The
post-commit ordering does not weaken the proximity guarantee because
(a) the merkle commit binds the prover to a specific polynomial
\emph{before} late points are sampled, and (b) the sumcheck and
STIR-query rounds operate over the same EqStatement that contains
both OOD and late claims, with constraint-batching $\gamma$ sampled
\emph{after} all claims are observed.

\paragraph{Multi-column variant.}  For $w$ columns, the per-column
adapter is invoked $w$ times with a shared challenger, with the
prover doing all-observes-then-all-opens (and the verifier mirroring).
This works but commits $w$ independent merkle trees, which is
wasteful at scale.  The Jagged + late-binding combined construction
in \S\ref{sec:design-jagged-late-binding} reduces this to one merkle
per shard.

\paragraph{Implementation.}
\texttt{crates/stark/src/whir\_late\_binding.rs}, built directly on
\texttt{p3-whir} internals (\texttt{CommitmentWriter},
\texttt{Verifier}, \texttt{EqStatement}); no upstream patch
required.  Tests: $24$ unit tests covering single-column, multi-
column, tampering rejection, and challenger-replay.

\subsection{Jagged + Late-Binding Reduction}
\label{sec:design-jagged-late-binding}
\label{sec:design-jagged-sumcheck}

The Jagged commitment (\S\ref{sec:design-jagged}) gives one merkle
per shard but commits $\widetilde{\mathbf{q}}$, not the per-chip
column MLEs that LogUp-GKR's leaf claims need.  The reduction:

\begin{construction}[Jagged sumcheck reduction]
\label{cons:jagged-sumcheck}
Inputs:
\begin{itemize}
  \item Dense polynomial $\widetilde{\mathbf{q}}$ committed via
        Construction~\ref{cons:late-binding}.
  \item Per-chip row coordinates $\mathbf{r}_{\mathrm{row},c}$ derived
        from GKR.
  \item Per-chip per-column row-MLE values
        $y_{c,j} = \widetilde{\mathbf{T}_c[:,j]}(\mathbf{r}_{\mathrm{row},c})$
        provided by the prover.
\end{itemize}

Steps:
\begin{enumerate}
  \item Prover observes all $y_{c,j}$ into the challenger.
  \item Verifier samples batching challenge $\gamma$.
  \item Prover and verifier compute initial claim
        $T = \sum_{c,j} \gamma^{\mathrm{idx}(c,j)} \cdot y_{c,j}$.
  \item Verifier defines the weight polynomial
        $\widetilde{w}(j) = \gamma^{\mathrm{idx}(c,j)} \cdot
        \mathsf{eq}(\mathrm{row}(j), \mathbf{r}_{\mathrm{row},c})$
        for $j$ in chip $c$ column-$j$'s range, $0$ elsewhere.
  \item Run sumcheck on $T = \sum_j \widetilde{\mathbf{q}}(j) \cdot
        \widetilde{w}(j)$ over $n$ variables (degree-$2$ rounds).
        After sumcheck, the verifier has
        $\widetilde{\mathbf{q}}(\mathbf{z}^*) \cdot
         \widetilde{w}(\mathbf{z}^*) = T_{\mathrm{final}}$
        where $\mathbf{z}^*$ is the sumcheck-derived point.
  \item Prover supplies $v^* = \widetilde{\mathbf{q}}(\mathbf{z}^*)$;
        verifier independently computes $\widetilde{w}(\mathbf{z}^*)$
        from the packing metadata, checks
        $v^* \cdot \widetilde{w}(\mathbf{z}^*) = T_{\mathrm{final}}$.
  \item $v^*$ is bound to the WHIR commit by registering
        $(\mathbf{z}^*, v^*)$ as a late claim
        (Construction~\ref{cons:late-binding}, step~4).
\end{enumerate}
\end{construction}

\begin{remark}[Convention reversal]
The sumcheck protocol's eval point is in last-variable-first order
$[r_n, r_{n-1}, \ldots, r_1]$ (because pair-consecutive folding
binds the last variable first), whereas
\texttt{p3\_multilinear\_util::Poly::eval\_base} expects MSB-first
$[x_0, x_1, \ldots, x_{n-1}]$.  The $\mathbf{z}^*$ passed to
\texttt{add\_late\_claim} is reversed; \texttt{verify\_jagged\_reduction}
mirrors this on the verifier side.  This was a real integration
bug found and fixed during implementation.
\end{remark}

\paragraph{Cost analysis.}  For a fibonacci shard ($N = 19$ chips,
$\sim 30$ avg cols, $n = 27$ packed):
\begin{itemize}
  \item Commit: $1$ merkle on $2^{27+r}$ codeword vs.
        $570$ merkles on $2^{17+r}$.
  \item Open: $1$ WHIR proof on $n=27$ vs. $570$ WHIR proofs on $n=17$.
  \item Sumcheck reduction: $n = 27$ rounds, degree-$2$
        polynomials, parallelisable per round.
  \item Per-shard proof bytes: $\sim 488$\,KB vs.~$\sim 29.5$\,MB
        ($60\times$ smaller).
\end{itemize}

Per-phase profile (parallelised, fibonacci shard):
\begin{tabular}{lr}
\toprule
WHIR commit on packed dense  & $1.7$\,s ($1.9\%$) \\
Sumcheck reduction (rayon)   & $0.16$\,s ($0.2\%$) \\
WHIR open phase              & $89.4$\,s ($97.9\%$) \\
\bottomrule
\end{tabular}

\paragraph{Implementation.}
\texttt{crates/stark/src/jagged\_late\_binding.rs}, $\sim 1200$ lines
including $11$ unit tests.

\subsection{Cross-Binding}
\label{sec:design-cross-binding}

\paragraph{Threat model.}  In the dual-commit configuration (FRI at
$\zeta$ + WHIR via Construction~\ref{cons:late-binding}), a
malicious prover could in principle commit to two different traces
$\mathbf{T}_A \neq \mathbf{T}_B$ --- the FRI side via the existing
\texttt{pcs.commit} on $\mathbf{T}_A$, the WHIR side via Jagged
late-binding on $\mathbf{T}_B$.  Both commits would individually
verify, but the proof would fail to attest to a single coherent
trace.

\begin{construction}[Brakedown-style cross-binding spot check]
\label{cons:cross-binding}
After both commits are observed:
\begin{enumerate}
  \item Sample $K$ random trace positions $i_1, \ldots, i_K$ from
        the post-both-commits challenger ($K = 80$ for $\sim 80$-bit
        soundness against single-position forgery).
  \item For each $i_j$, the prover provides:
        \begin{enumerate}
          \item The trace value $\mathbf{T}[i_j]$.
          \item Merkle paths against the FRI MMCS at the codeword
                position(s) corresponding to $i_j$.
          \item Merkle paths against the WHIR MMCS at the codeword
                position corresponding to $i_j$.
        \end{enumerate}
  \item Verifier checks both paths via \texttt{Mmcs::verify\_batch}
        and that both reveal the same value.
\end{enumerate}
\end{construction}

\paragraph{Soundness.}  If $\mathbf{T}_A \neq \mathbf{T}_B$, they
differ at $\geq 1$ position; $K$ random positions catch the
divergence with probability $1 - (1 - 1/N)^K \approx 1 - 2^{-K}$
for $N = 2^{17}, K = 80$.

\paragraph{Status.}  Both sides fully implemented at host level:
prover populates FRI- and WHIR-side merkle paths
(\texttt{populate\_fri\_paths},
\texttt{prove\_cross\_binding\_spot\_check}); verifier checks both
via \texttt{Mmcs::verify\_batch}
(\texttt{verify\_cross\_binding\_spot\_check\_fri},
\texttt{verify\_cross\_binding\_spot\_check\_whir}).  The
verifier dispatcher
\texttt{verify\_cross\_binding\_for\_kb} extracts the WHIR merkle
root from the jagged bundle's \texttt{WhirProof.initial\_commitment}
and reconstructs the codeword shape via the public helper
\texttt{whir\_codeword\_dims(log\_dense\_size)}, which mirrors
\texttt{make\_chip\_late\_binding\_pcs}'s folding-factor derivation.
Tamper rejection at the merkle layer is validated by
\texttt{cross\_binding\_whir\_rejects\_tampered\_path}.

\paragraph{Implementation.}
Helpers in
\texttt{crates/stark/src/jagged\_late\_binding.rs}
(\texttt{populate\_fri\_paths},
\texttt{verify\_cross\_binding\_spot\_check\_fri},
\texttt{verify\_cross\_binding\_spot\_check\_whir},
\texttt{prove\_cross\_binding\_spot\_check}); the codeword-dimension
helper \texttt{whir\_codeword\_dims} in
\texttt{crates/stark/src/whir\_late\_binding.rs}; prover wiring in
\texttt{crates/stark/src/prover.rs}
(\texttt{try\_compute\_cross\_binding\_fri},
\texttt{prove\_jagged\_late\_binding\_with\_cross\_binding});
verifier wiring in \texttt{crates/stark/src/verifier.rs}
(\texttt{try\_verify\_cross\_binding},
\texttt{verify\_cross\_binding\_for\_kb}).

\subsection{Soundness Layers Summary}

The combined pipeline composes the following soundness layers:

\begin{tabular}{@{}lll@{}}
\toprule
\textbf{Layer}                       & \textbf{Mechanism}            & \textbf{Soundness Bound} \\
\midrule
Transition constraints               & Zerocheck (sumcheck)          & $d \cdot m / |\mathbb{F}_q|$ \\
Lookup arguments                     & LogUp-GKR + per-layer sumcheck & $3 \cdot \log_2 N / |\mathbb{F}_q|$ \\
Row-MLE binding to commit            & Late-binding WHIR             & WHIR proximity \\
FRI $\leftrightarrow$ WHIR coherence  & Cross-binding spot-check (C.1) & $1 - (1 - 1/N)^K$ \\
\bottomrule
\end{tabular}

For our parameters ($\mathbb{F}_q \approx 2^{124}$, $m \leq 27$,
$d = 2$, $N \leq 2^{27}$, $K = 80$), each layer's soundness error is
$\leq 2^{-80}$.

\subsection{Implementation Activation Matrix}

\begin{tabular}{@{}ll@{}}
\toprule
\textbf{Configuration}                                  & \textbf{Env flags} \\
\midrule
FRI (default, unchanged)                                & \emph{(none)} \\
WHIR fast path + per-chip late-binding                  & \texttt{ZIREN\_USE\_WHIR=1} \\
WHIR fast path + jagged late-binding                    & \texttt{ZIREN\_USE\_WHIR=1
                                                            ZIREN\_LATE\_BINDING=jagged} \\
+ cross-binding (FRI + WHIR sides)                      & + \texttt{ZIREN\_CROSS\_BINDING=1} \\
\bottomrule
\end{tabular}

Production deployment should use the fourth row (jagged late-binding
with cross-binding).  End-to-end smoke tests on fibonacci, JSON,
keccak, and large-sum (multi-shard) workloads validate the first
three rows; the fourth row's merkle-layer correctness is validated
by lib-level round-trip and tamper-rejection unit tests against
the same MMCS instance the prover uses.
